\documentclass[11pt]{article}

\usepackage[a4paper,margin=1in]{geometry}
\usepackage[T1]{fontenc}
\usepackage[utf8]{inputenc}
\usepackage{amsmath}
\usepackage{graphicx}
\usepackage{booktabs}
\usepackage{tabularx}
\usepackage{array}
\usepackage{subcaption}
\usepackage{float}
\usepackage{hyperref}
\usepackage{enumitem}

\hypersetup{
    colorlinks=true,
    linkcolor=blue,
    citecolor=blue,
    urlcolor=blue
}

\graphicspath{{figures/scheduling/}}

\title{Finite-Blocklength Latency Minimization for Uplink and Downlink Communications:\\Per-User Scheduling Results}
\author{}
\date{July 6, 2026}

\begin{document}

\maketitle

\begin{abstract}
This document focuses only on per-user scheduling behavior. Throughout, the
\emph{Training Only Method} is used as the benchmark for the
\emph{Training + Testing Method}. The clearest scheduling gap appears in the
downlink. In payload completion, the benchmark finishes with block counts
\((4,3,3)\), whereas the \emph{Training + Testing Method} expands that to
\((9,3,7)\). In continuous communication, both downlink methods keep five
scheduled blocks per user, but the benchmark still preserves much stronger
service on users 1 and 2. In the uplink, the scheduling structures are already
close, and the main visible difference is limited to a single extra payload
block for user 1 under the \emph{Training + Testing Method}.
\end{abstract}

\section{Scope and Naming}

This document is intentionally schedule-only. The terminology is aligned with
the companion theory report:
\begin{itemize}[leftmargin=1.5em]
    \item \textbf{Training Only Method}: direct online optimization during scheduling,
    \item \textbf{Training + Testing Method}: a separate training stage followed by testing-stage inference,
    \item \textbf{Payload completion}: the scenario where each user continues until its remaining payload is exhausted, and
    \item \textbf{Continuous communication}: the scenario with a prescribed communication horizon and prescribed blockwise demand sequence.
\end{itemize}

The \emph{Training Only Method} is treated as the benchmark, and the
\emph{Training + Testing Method} is evaluated against it.

\section{Downlink}

The downlink schedule results separate the two methods most clearly. The main
difference is not only how long the schedule is, but how efficiently service is
distributed across users once multiuser coupling becomes difficult.

\subsection{Payload Completion Results and Comparison}

Downlink payload completion shows the strongest benchmark advantage. The
\emph{Training Only Method} finishes with user block counts \((4,3,3)\), while
the \emph{Training + Testing Method} expands the same task to \((9,3,7)\).
That means the main schedule tail is carried by users 0 and 2.

\subsubsection*{Training Only Method}

\begin{itemize}[leftmargin=1.5em,itemsep=0.2em,topsep=0.3em]
    \item \textbf{User 0}: 4 blocks, total \(n=198\), served \(1000/1000\), with 3 skipped blocks.
    \item \textbf{User 1}: 3 blocks, total \(n=140\), served \(2000/2000\), with no skipped blocks.
    \item \textbf{User 2}: 3 blocks, total \(n=199\), served \(3000/3000\), with no skipped blocks.
\end{itemize}

\subsubsection*{Training + Testing Method}

\begin{itemize}[leftmargin=1.5em,itemsep=0.2em,topsep=0.3em]
    \item \textbf{User 0}: 9 blocks, total \(n=427\), served \(1000/1000\), with 7 skipped blocks.
    \item \textbf{User 1}: 3 blocks, total \(n=164\), served \(2000/2000\), with no skipped blocks.
    \item \textbf{User 2}: 7 blocks, total \(n=512\), served \(3000/3000\), with 3 skipped blocks.
\end{itemize}

Figure~\ref{fig:dl-payload-schedule} makes the same point visually. The
benchmark remains compact, with average chosen sub-blocklengths approximately
\((49.5,46.7,66.3)\). The \emph{Training + Testing Method} pushes users 1 and 2
to heavier average sub-blocklengths, approximately \((47.4,54.7,73.1)\), while
also fragmenting users 0 and 2 across more blocks. So the remaining gap is not
just a small timing shift; it is a materially longer and less synchronized
schedule.

\begin{figure}[H]
\centering
\begin{subfigure}[t]{0.48\textwidth}
    \centering
    \includegraphics[width=\textwidth,height=0.27\textheight,keepaspectratio]{dl_payload_training_only.png}
    \caption{Training Only Method}
\end{subfigure}
\hfill
\begin{subfigure}[t]{0.48\textwidth}
    \centering
    \includegraphics[width=\textwidth,height=0.27\textheight,keepaspectratio]{dl_payload_training_testing.png}
    \caption{Training + Testing Method}
\end{subfigure}
\caption{Downlink payload-completion schedule comparison. The benchmark keeps the schedule compact, while the \emph{Training + Testing Method} stretches users 0 and 2 over a much longer block tail.}
\label{fig:dl-payload-schedule}
\end{figure}

\clearpage
\subsection{Continuous Communication Results and Comparison}

In downlink continuous communication, both methods still allocate five scheduled
blocks per user, so the difference is no longer block count. The important
question becomes which users actually receive usable service within that fixed
horizon.

\subsubsection*{Training Only Method}

\begin{itemize}[leftmargin=1.5em,itemsep=0.2em,topsep=0.3em]
    \item \textbf{User 0}: 5 blocks, total \(n=250\), served \(0/1000\), with 5 zero-service blocks.
    \item \textbf{User 1}: 5 blocks, total \(n=140\), served \(2000/2000\), and fully met the target.
    \item \textbf{User 2}: 5 blocks, total \(n=264\), served \(2400/3000\), with 1 zero-service block and 600 unserved bits.
\end{itemize}

\subsubsection*{Training + Testing Method}

\begin{itemize}[leftmargin=1.5em,itemsep=0.2em,topsep=0.3em]
    \item \textbf{User 0}: 5 blocks, total \(n=250\), served \(0/1000\), with 5 zero-service blocks.
    \item \textbf{User 1}: 5 blocks, total \(n=209\), served \(1928/2000\), with 1 partial block and 72 unserved bits.
    \item \textbf{User 2}: 5 blocks, total \(n=400\), served \(0/3000\), with 5 zero-service blocks.
\end{itemize}

The benchmark therefore uses the same five-block horizon more effectively.
For the \emph{Training Only Method}, the average chosen sub-blocklengths are
approximately \((50.0,28.0,52.8)\). For the \emph{Training + Testing Method},
they increase to approximately \((50.0,41.8,80.0)\). User 2 is the clearest
contrast: the \emph{Training + Testing Method} spends the full average
blocklength budget on that user but still delivers no service. The limitation
here is therefore feasibility within the fixed horizon rather than lack of
scheduling opportunities.

\begin{figure}[H]
\centering
\begin{subfigure}[t]{0.48\textwidth}
    \centering
    \includegraphics[width=\textwidth,height=0.27\textheight,keepaspectratio]{dl_continuous_training_only.png}
    \caption{Training Only Method}
\end{subfigure}
\hfill
\begin{subfigure}[t]{0.48\textwidth}
    \centering
    \includegraphics[width=\textwidth,height=0.27\textheight,keepaspectratio]{dl_continuous_training_testing.png}
    \caption{Training + Testing Method}
\end{subfigure}
\caption{Downlink continuous-communication schedule comparison. Both methods retain five scheduled blocks per user, but only the benchmark converts that horizon into stronger service on users 1 and 2.}
\label{fig:dl-continuous-schedule}
\end{figure}

\clearpage
\section{Uplink}

The uplink results are much closer to the benchmark. In both scenarios, the
\emph{Training + Testing Method} largely preserves the same per-user schedule
shape and differs mainly through small blocklength redistributions.

\subsection{Payload Completion Results and Comparison}

Uplink payload completion keeps the same overall front-loaded structure in both
methods: most service is delivered early, and the last block is used to clear a
small remainder. The benchmark finishes with block counts \((2,2,3)\), while
the \emph{Training + Testing Method} gives \((2,3,3)\). The only structural
difference is the extra block for user 1.

\subsubsection*{Training Only Method}

\begin{itemize}[leftmargin=1.5em,itemsep=0.2em,topsep=0.3em]
    \item \textbf{User 0}: 2 blocks, total \(n=68\), served \(1000/1000\), and cleared in 2 blocks.
    \item \textbf{User 1}: 2 blocks, total \(n=119\), served \(2000/2000\), and cleared in 2 blocks.
    \item \textbf{User 2}: 3 blocks, total \(n=183\), served \(3000/3000\), with a small final clean-up block.
\end{itemize}

\subsubsection*{Training + Testing Method}

\begin{itemize}[leftmargin=1.5em,itemsep=0.2em,topsep=0.3em]
    \item \textbf{User 0}: 2 blocks, total \(n=67\), served \(1000/1000\), and cleared in 2 blocks.
    \item \textbf{User 1}: 3 blocks, total \(n=129\), served \(2000/2000\), with one extra block versus the benchmark.
    \item \textbf{User 2}: 3 blocks, total \(n=182\), served \(3000/3000\), with a small final clean-up block.
\end{itemize}

So the benchmark edge is very localized. Users 0 and 2 are essentially matched,
while user 1 is the only case where the \emph{Training + Testing Method}
extends the schedule. Figure~\ref{fig:ul-payload-schedule} shows that both
methods still keep the same basic policy of a large early allocation followed
by a sharply reduced final block.

\begin{figure}[H]
\centering
\begin{subfigure}[t]{0.48\textwidth}
    \centering
    \includegraphics[width=\textwidth,height=0.27\textheight,keepaspectratio]{ul_payload_training_only.png}
    \caption{Training Only Method}
\end{subfigure}
\hfill
\begin{subfigure}[t]{0.48\textwidth}
    \centering
    \includegraphics[width=\textwidth,height=0.27\textheight,keepaspectratio]{ul_payload_training_testing.png}
    \caption{Training + Testing Method}
\end{subfigure}
\caption{Uplink payload-completion schedule comparison. The two methods follow almost the same structure, with the only material change being one extra block for user 1 under the \emph{Training + Testing Method}.}
\label{fig:ul-payload-schedule}
\end{figure}

\clearpage
\subsection{Continuous Communication Results and Comparison}

Uplink continuous communication is the closest schedule match in the document.
Both methods keep exactly five scheduled blocks for every user, and all target
payloads are fully served.

\subsubsection*{Training Only Method}

\begin{itemize}[leftmargin=1.5em,itemsep=0.2em,topsep=0.3em]
    \item \textbf{User 0}: 5 blocks, total \(n=79\), served \(1000/1000\), and fully served in the 5-block horizon.
    \item \textbf{User 1}: 5 blocks, total \(n=130\), served \(2000/2000\), and fully served in the 5-block horizon.
    \item \textbf{User 2}: 5 blocks, total \(n=201\), served \(3000/3000\), and fully served in the 5-block horizon.
\end{itemize}

\subsubsection*{Training + Testing Method}

\begin{itemize}[leftmargin=1.5em,itemsep=0.2em,topsep=0.3em]
    \item \textbf{User 0}: 5 blocks, total \(n=77\), served \(1000/1000\), and is slightly shorter than the benchmark.
    \item \textbf{User 1}: 5 blocks, total \(n=137\), served \(2000/2000\), and is slightly longer than the benchmark.
    \item \textbf{User 2}: 5 blocks, total \(n=196\), served \(3000/3000\), and is slightly shorter than the benchmark.
\end{itemize}

This is therefore not a different schedule topology problem. It is a small
reallocation of sub-blocklength inside the same five-block structure. Relative
to the benchmark, the \emph{Training + Testing Method} shortens users 0 and 2
slightly and lengthens user 1 slightly, but the overall shape stays the same.

\begin{figure}[H]
\centering
\begin{subfigure}[t]{0.48\textwidth}
    \centering
    \includegraphics[width=\textwidth,height=0.27\textheight,keepaspectratio]{ul_continuous_training_only.png}
    \caption{Training Only Method}
\end{subfigure}
\hfill
\begin{subfigure}[t]{0.48\textwidth}
    \centering
    \includegraphics[width=\textwidth,height=0.27\textheight,keepaspectratio]{ul_continuous_training_testing.png}
    \caption{Training + Testing Method}
\end{subfigure}
\caption{Uplink continuous-communication schedule comparison. The five-block per-user structure is preserved exactly, and the remaining difference is only a mild redistribution of chosen sub-blocklength across users.}
\label{fig:ul-continuous-schedule}
\end{figure}

\section{Schedule Takeaways}

The scheduling results are most naturally read benchmark-first. In downlink
payload completion, the main shortfall of the \emph{Training + Testing Method}
is schedule fragmentation, especially on users 0 and 2. In downlink continuous
communication, the main issue is even stronger: the horizon is present, but the
resulting service remains weak, especially for user 2.

In the uplink, the schedule gap is already small. Payload completion differs
mainly through one extra user-1 block, and continuous communication preserves
the same five-block structure for all users. So the present scheduling gap is
mostly a downlink issue, while the uplink is already close to the benchmark at
the per-user schedule level.

\end{document}
